\documentclass[11pt]{article}
\usepackage[T1]{fontenc}
\usepackage[utf8]{inputenc}
\usepackage{lmodern}
\usepackage[margin=1in]{geometry}
\usepackage{array}
\usepackage{booktabs}
\usepackage{enumitem}
\usepackage{longtable}
\usepackage{tabularx}
\PassOptionsToPackage{hyphens}{url}
\usepackage[hidelinks]{hyperref}
\usepackage{xurl}
\setlength{\parindent}{0pt}
\setlength{\parskip}{0.6em}
\setlist[itemize]{leftmargin=1.4em}
\setlist[enumerate]{leftmargin=1.6em}
\emergencystretch=2em
\sloppy

\title{Applying ISO/IEC/IEEE 42010:2022\\to an Enterprise Architecture Description}
\author{}
\date{}

\begin{document}
\maketitle

\section*{Purpose}
This guide explains how to apply ISO/IEC/IEEE 42010:2022 to a real enterprise architecture description (AD). It is project-facing rather than theoretical: the goal is to help an architecture team produce an enterprise AD that is reviewable, traceable, and fit for governance.

\section*{What ISO 42010 Gives You}
ISO/IEC/IEEE 42010:2022 gives you requirements for the structure and expression of an architecture description, not a mandatory enterprise architecture method. The official ISO record states that the standard covers architecture descriptions for software, systems, enterprises, systems of systems, product lines, service lines, technologies, and business domains, while leaving process, tools, methods, notations, and media open \cite{iso42010}.

For enterprise architecture work, that means ISO 42010 is best used as a \textbf{discipline for organizing architectural evidence}. It helps you show:
\begin{itemize}
    \item what enterprise or scope segment is being described,
    \item which stakeholders and concerns matter,
    \item which viewpoints and views address those concerns,
    \item which modeling conventions are in use,
    \item and what decisions, rationale, and correspondences make the description coherent \cite{iso42010,cm42010}.
\end{itemize}

\section*{Important Constraint}
The workflow below is an \textbf{application pattern}, not an ISO-prescribed lifecycle. ISO 42010 does not tell you to use TOGAF, ArchiMate, BPMN, UML, capability maps, value streams, or repository tools. It tells you to make your choices explicit and tie them to stakeholder concerns \cite{iso42010}.

\section*{A Practical Application Sequence}

\subsection*{1. Define the entity of interest and scope}
Start by naming the enterprise subject of the architecture description. In enterprise work, the entity of interest might be:
\begin{itemize}
    \item the whole enterprise,
    \item a business unit,
    \item a platform,
    \item an operating model,
    \item a transformation program,
    \item or a cross-functional capability such as identity, finance, customer support, or cloud governance.
\end{itemize}

Document the boundary explicitly:
\begin{itemize}
    \item what is in scope,
    \item what is out of scope,
    \item which time horizon the AD addresses,
    \item and which external actors or environmental forces matter.
\end{itemize}

This aligns with the standard's distinction between the entity of interest and the AD that expresses its architecture \cite{iso42010}.

\subsection*{2. Identify stakeholders and concerns}
Before producing views, capture who the AD must serve and what each audience needs answered. In enterprise architecture, common stakeholders include:
\begin{itemize}
    \item executive sponsors,
    \item business owners,
    \item enterprise architects,
    \item domain architects,
    \item security and risk leaders,
    \item delivery teams,
    \item operations teams,
    \item compliance and audit,
    \item data governance,
    \item and external partners or regulators.
\end{itemize}

Typical concerns include strategic alignment, investment prioritization, business capability coverage, security posture, compliance obligations, interoperability, scalability, data ownership, operational resilience, and migration risk. In ISO 42010 terms, stakeholder concerns are the reason viewpoints and views exist \cite{cm42010}.

\subsection*{3. Select the minimum set of viewpoints}
Define viewpoints that can address the identified concerns. A viewpoint is not just a diagram name; it is a specification for how a view will be constructed. The 2022 framing allows one viewpoint to govern one or more views \cite{nora42010}.

A practical enterprise architecture viewpoint set may include:
\begin{itemize}
    \item \textbf{Context viewpoint} for scope, boundaries, and external actors,
    \item \textbf{Business viewpoint} for capabilities, value streams, organization, and processes,
    \item \textbf{Information viewpoint} for data domains, ownership, semantics, and flows,
    \item \textbf{Application viewpoint} for systems, services, interfaces, and dependencies,
    \item \textbf{Technology viewpoint} for platforms, environments, hosting, and infrastructure,
    \item \textbf{Security and risk viewpoint} for controls, trust boundaries, threats, and obligations,
    \item \textbf{Transition viewpoint} for roadmap states, work packages, dependencies, and migration steps.
\end{itemize}

These are example viewpoint categories, not ISO-mandated names.

\subsection*{4. Define model kinds and notation rules}
The ISO abstract specifically names model kinds as part of the support structure for an AD \cite{iso42010}. In enterprise practice, this means you should state which model kinds and notations are allowed within each viewpoint.

Example:
\begin{itemize}
    \item business capability map,
    \item organization interaction model,
    \item application interaction model,
    \item information flow model,
    \item deployment model,
    \item control mapping matrix,
    \item transition roadmap.
\end{itemize}

For each model kind, define:
\begin{itemize}
    \item allowed notation or notation family,
    \item required element types,
    \item naming rules,
    \item level of abstraction,
    \item mandatory metadata,
    \item and review criteria.
\end{itemize}

\subsection*{5. Produce architecture views that answer real questions}
Do not create views because a framework template says they should exist. Create them because they resolve stakeholder concerns. Every view in the AD should be traceable to at least one concern.

A useful enterprise view should answer a concrete question, such as:
\begin{itemize}
    \item Which business capabilities are in scope for the transformation?
    \item Which applications support each critical capability?
    \item Which systems exchange regulated data?
    \item Where are the control points and trust boundaries?
    \item What is the current state, target state, and transition path?
\end{itemize}

\subsection*{6. Record architecture decisions and rationale}
An enterprise AD is incomplete if it contains diagrams without decisions. Record major decisions such as:
\begin{itemize}
    \item platform standardization choices,
    \item sourcing choices,
    \item integration patterns,
    \item data ownership allocations,
    \item security boundary choices,
    \item governance model choices,
    \item target-state simplification or consolidation decisions.
\end{itemize}

For each decision, record:
\begin{itemize}
    \item the decision statement,
    \item the options considered,
    \item the selected option,
    \item the rationale,
    \item assumptions,
    \item constraints,
    \item consequences,
    \item and related views or models.
\end{itemize}

This preserves traceability and makes reviews materially stronger \cite{cm42010}.

\subsection*{7. Define correspondences across views}
Correspondence is one of the most valuable ISO 42010 ideas for enterprise architecture. It allows you to state how elements in different views relate and how consistency is checked \cite{cm42010}.

Examples of enterprise correspondences:
\begin{itemize}
    \item business capability to application support mapping,
    \item application to information domain ownership mapping,
    \item control objective to technology enforcement point mapping,
    \item transition work package to affected capability or system mapping,
    \item target-state element to current-state predecessor mapping.
\end{itemize}

Without correspondences, enterprise architecture often becomes a disconnected set of diagrams.

\subsection*{8. Review the AD against concerns, not aesthetics}
A reviewable enterprise AD should be assessed by asking:
\begin{itemize}
    \item Are the right stakeholders identified?
    \item Are their concerns explicit?
    \item Does every significant concern have an addressing viewpoint and one or more views?
    \item Are model kinds and notation rules defined?
    \item Are major decisions and rationale captured?
    \item Are cross-view correspondences explicit?
    \item Is the scope and boundary of the entity of interest clear?
\end{itemize}

This is much closer to the intent of ISO 42010 than simply checking whether a document contains many diagrams.

\subsection*{9. Maintain the architecture description as a governed artifact}
An enterprise AD is not a one-time slide deck. Treat it as a controlled work product with:
\begin{itemize}
    \item versioning,
    \item ownership,
    \item review cadence,
    \item traceability to decisions,
    \item traceability to programs or portfolios,
    \item and change control for key views and correspondences.
\end{itemize}

\section*{A Minimal Enterprise Architecture Description Structure}
The following structure works well as a practical minimum:

\begin{center}
\renewcommand{\arraystretch}{1.2}
\begin{tabularx}{\textwidth}{>{\raggedright\arraybackslash}p{0.27\textwidth} X}
\toprule
\textbf{Section} & \textbf{What to include} \\
\midrule
Purpose and scope & Why the AD exists, entity of interest, boundaries, assumptions, time horizon. \\
Stakeholders and concerns & Named audiences and the concerns the AD must address. \\
Viewpoint catalogue & Defined viewpoints, intended concerns, allowed model kinds, notation conventions. \\
Architecture views & The actual views and models used to answer stakeholder questions. \\
Decisions and rationale log & Major decisions, alternatives, rationale, consequences, dependencies. \\
Correspondence set & Cross-view mappings, consistency rules, traceability relationships. \\
Review and governance & Ownership, status, review dates, approval path, maintenance expectations. \\
\bottomrule
\end{tabularx}
\end{center}

\section*{Example Stakeholder-to-View Mapping}
\begin{center}
\renewcommand{\arraystretch}{1.2}
\begin{tabularx}{\textwidth}{>{\raggedright\arraybackslash}p{0.18\textwidth} >{\raggedright\arraybackslash}p{0.25\textwidth} >{\raggedright\arraybackslash}p{0.22\textwidth} X}
\toprule
\textbf{Stakeholder} & \textbf{Concern} & \textbf{Viewpoint} & \textbf{Example view} \\
\midrule
Executive sponsor & Investment alignment and business outcome visibility & Business / transition viewpoint & Capability heatmap with phased roadmap \\
Security lead & Control coverage and trust boundaries & Security and risk viewpoint & Trust boundary and control enforcement view \\
Application owner & Integration dependencies and system impact & Application viewpoint & Application interaction map \\
Data governance lead & Ownership, lineage, and regulated data movement & Information viewpoint & Information domain and flow view \\
Operations lead & Runtime dependencies and resilience & Technology viewpoint & Platform dependency and resilience view \\
\bottomrule
\end{tabularx}
\end{center}

\section*{Common Failure Modes}
Common enterprise architecture problems that ISO 42010 helps prevent include:
\begin{itemize}
    \item diagrams with no identified stakeholder or concern,
    \item large document sets with no explicit viewpoint logic,
    \item inconsistent terminology across business, application, and technology models,
    \item target-state claims with no rationale,
    \item roadmap views that are disconnected from architecture decisions,
    \item and review cycles that focus on diagram appearance rather than concern coverage and traceability.
\end{itemize}

\section*{Bottom Line}
Apply ISO/IEC/IEEE 42010:2022 to enterprise architecture by treating the architecture description as a governed, concern-driven body of evidence. Start with the entity of interest, identify stakeholders and concerns, define viewpoints and model kinds, build views that answer specific questions, record decisions and rationale, and maintain explicit correspondences across the description. That is the shortest path to an enterprise AD that is useful for governance, communication, and change execution \cite{iso42010,cm42010}.

\section*{References}
\begin{thebibliography}{9}
\bibitem{iso42010}
International Organization for Standardization (ISO), \emph{ISO/IEC/IEEE 42010:2022 - Software, systems and enterprise --- Architecture description}, official standard record and abstract. Available at: \url{https://www.iso.org/standard/74393.html}

\bibitem{cm42010}
ISO/IEC/IEEE 42010 community site, \emph{ISO/IEC/IEEE 42010: Conceptual Model}. Available at: \url{https://www.iso-architecture.org/42010/cm/}

\bibitem{nora42010}
NORA Online, \emph{ISO 42010}. Used here only as a public summary of terminology and framing differences relevant to the 2022 edition. Available at: \url{https://www.noraonline.nl/wiki/ISO_42010}
\end{thebibliography}

\end{document}



